\documentclass[11pt]{article}
\usepackage[margin=1in]{geometry}
\usepackage[T1]{fontenc}\usepackage[utf8]{inputenc}
\usepackage{lmodern,amsmath,amssymb,amsthm,mathtools,booktabs,microtype}
\usepackage[colorlinks=true,linkcolor=blue,citecolor=blue,urlcolor=blue]{hyperref}
\hypersetup{pdftitle={Polynomial-depth Bayesian learning with continuous neuron gates},pdfauthor={ChatGPT-6 Astra; Claude Opus 5.5; Kaarel Hänni; Dmitry Vaintrob}}
\newtheorem{theorem}{Theorem}[section]\newtheorem{lemma}[theorem]{Lemma}
\newtheorem{corollary}[theorem]{Corollary}\newtheorem{proposition}[theorem]{Proposition}
\theoremstyle{remark}\newtheorem{remark}[theorem]{Remark}
\newcommand{\E}{\mathbb E}\newcommand{\R}{\mathbb R}\newcommand{\Pp}{\mathbb P}
\newcommand{\KL}{\operatorname{KL}}\newcommand{\whR}{\widehat R}
\title{Polynomial-depth Bayesian learning\\with continuous neuron gates}
\author{ChatGPT-6 Astra \and Claude Opus 5.5 \and Kaarel H\"anni \and
Dmitry Vaintrob\thanks{Authors are listed in alphabetical order.}}
\date{17 November 2026}
\begin{document}\maketitle
\begin{abstract}
We give a Bayesian neural-network prior whose sample complexity is polynomial jointly in the width and depth of a comparison network, uniformly over every larger finite learner width. The architecture is fixed. Each hidden neuron has a positive gamma gate shared by its outgoing weights, and the remaining edge factors have independent Student distributions. All actual weights are continuous and nonzero almost surely. This hierarchical group-shrinkage prior is a substantive change from independent Gaussian weights; no claim is made for the latter. For any globally Lipschitz hidden activation, biases at every layer, and a bounded Lipschitz output link, the ordinary posterior mean under a fixed Gaussian likelihood has expected risk at most comparison risk plus a nearly $1/n$ remainder. For a width-$C$, depth-$H$ comparator with bounded coefficients, the remainder is, up to logarithms, $O((dC+HC^2)H/n)$, independent of learner width. Coefficient magnitudes enter only logarithmically. For Boolean circuits, restoring gate implementations remove the extra depth factor and give a remainder of order $\widetilde O((dC+HC^2)/n)$. For ReLU and tanh, extra learner depth can be carried through one scalar channel at a nearly linear additive cost. A high-probability companion separates the noise floor and has an arbitrarily small multiplicative slack on approximation error. Both allow conditional sub-Gaussian label noise. The proof suppresses the total gate mass outside an embedded computation; the number of possible neuron placements cancels the probability cost of all active gates, even when intermediate computations are shared. No robustness restriction is imposed on the prior or on the comparison network. Conversely, typical prior draws compress to networks with polynomial depth and description complexity; this gives an actual function-prior complexity converse, not only a learning upper bound.
\end{abstract}

\section{A fixed architecture with neuron-level shrinkage}
A useful neural predictor may be much smaller than the architecture used for Bayesian learning. We seek an oracle inequality comparing the learned function with such a smaller network, with a cost polynomial in both its width and its depth. Sharing intermediate computations matters: replacing a general network by a computation tree can increase its size exponentially in depth.

We use one continuous gate per hidden neuron, shared by all its outgoing weights. A gate near zero suppresses that neuron's contribution to every later use. A gate of order one permits the computation to be reused without a separate rare-event cost at each outgoing edge. The gate distribution has shape inversely proportional to learner width, so the total gate mass of a layer has a width-independent distribution.

This is a hierarchical group-sparsity prior, related in spirit to automatic relevance determination. It samples neither an architecture nor a discrete network size, and the finite network remains dense almost surely. Nevertheless, the dependence among outgoing weights is essential. The theorem does not settle the same question for independent Gaussian coordinates, or for an arbitrary conventional prior. It concerns exact Bayesian prediction, not the computational cost of evaluating the posterior.

\section{The model}
Fix hidden depth $H\ge1$, learner width $N$, and input domain
$\mathcal X\subseteq[-1,1]^d$. Let $\phi:\R\to\R$ be globally $A$-Lipschitz, with $A\ge1$ and $b=|\phi(0)|$. Let $\psi:\R\to[-1,1]$ be $L_\psi$-Lipschitz, with $L_\psi\ge1$. ReLU, GELU, tanh, sigmoid, softplus and SiLU are examples of allowed hidden activations. The output link can be tanh or clipping.

Fix $0<\kappa\le1$ and $\nu>2$. Independently sample gates and factors
\[
 S_j^\ell\sim\operatorname{Gamma}(\kappa/N,1),
 \qquad G\sim t_\nu,
\]
where the second expression denotes the symmetric Student distribution with density proportional to $(1+t^2/\nu)^{-(\nu+1)/2}$. Every edge factor and every bias factor has this same law. The network is
\begin{align}
 h_i^1(x)&=\phi\!\left(G^1_{i0}+\sum_{j=1}^dG^1_{ij}x_j\right),\\
 h_i^\ell(x)&=\phi\!\left(G^\ell_{i0}
       +\sum_{j=1}^N S_j^{\ell-1}G^\ell_{ij}h_j^{\ell-1}(x)\right),
       &&2\le\ell\le H,\\
 z_w(x)&=G^{\rm out}_0+\sum_{j=1}^NS_j^HG^{\rm out}_jh_j^H(x),
 &f_w(x)&=\psi(z_w(x)).
 \label{eq:model}
\end{align}
Write $\mu_N$ for the prior on all latent gates and factors. Its pushforward to actual weights is also the prior used by the learner. Actual weights leaving one neuron share its gate and are dependent. The coordinate vector has full support and is absolutely continuous; all coordinates are nonzero almost surely. The total gate mass at each layer is exactly $\operatorname{Gamma}(\kappa,1)$, for every $N$.

\begin{lemma}[ReLU equivalence to incoming-row group shrinkage]\label{lem:row-gates}
For ReLU, the same function prior is obtained by multiplying each neuron's entire incoming row and its bias by its own gate, and leaving the final readout factors ungated. In this representation,
\[
 \widetilde h_i^1=S_i^1\operatorname{ReLU}\!\left(G^1_{i0}+\sum_jG^1_{ij}x_j\right),
 \qquad
 \widetilde h_i^\ell=
 \operatorname{ReLU}\!\left(S_i^\ell G^\ell_{i0}
       +\sum_jS_i^\ell G^\ell_{ij}\widetilde h_j^{\ell-1}\right),
\]
and $z=G^{\rm out}_0+\sum_jG^{\rm out}_j\widetilde h_j^H$.
\end{lemma}
\begin{proof}
Positive homogeneity and induction give $\widetilde h_i^\ell=S_i^\ell h_i^\ell$, where $h$ is the feature in~\eqref{eq:model}. Substitution in the readout proves exact equality for every parameter realization.
\end{proof}
Thus, for ReLU, the construction can be viewed as ordinary per-neuron incoming-group shrinkage, including biases. This identity does not apply to tanh; its outgoing-gate representation remains the model definition.

A comparator $g$ has the same depth, activation and output link, hidden-layer widths $m_\ell\le m\le N$, and all its weights and biases bounded in magnitude by $B\ge1$. It may share intermediate computations arbitrarily. Define
\begin{align}
 V&=\sum_{\ell=1}^H m_\ell,\\
 P&=(d+1)m_1+\sum_{\ell=2}^H m_\ell(m_{\ell-1}+1)+m_H+1,\\
 Q&=P+V+H.
 \label{eq:counts}
\end{align}
Here $P$ counts the full blocks between selected teacher neurons, including coefficients equal to zero. The result is a width-based bound, rather than a bound only in the number of nonzero edges.

\section{The learning guarantees}
Suppose the observations are IID, with regression mean
$f_0(x)=\E[Y\mid X=x]\in[-1,1]$. Assume $\xi=Y-f_0(X)$ is conditionally $\sigma$-sub-Gaussian:
\[
 \E[e^{t\xi}\mid X]\le e^{\sigma^2t^2/2}\qquad(t\in\R).
\]
The input law is otherwise arbitrary; observed labels may be unbounded. Set
\[
 R(f)=\E(f(X)-Y)^2,\quad R_0=\E\xi^2,\quad s=1+\sigma^2.
\]
For a fixed inverse temperature $\beta>0$, use the Gaussian-likelihood posterior
\[
 d\Pi_{S_n,\beta}(w)=Z_{S_n,\beta}^{-1}
 e^{-\beta\sum_{i=1}^n(f_w(X_i)-Y_i)^2}\,d\mu_N(w),
 \qquad \bar f_{S_n,\beta}=\int f_w\,d\Pi_{S_n,\beta}.
\]
The observation variance is $1/(2\beta)$. The predictor is the ordinary final posterior mean trained on all observations.

Constants $C$ below depend only on the fixed gate and factor distributions. Put $\mu=(\E G^2)^{1/2}$ and choose
\begin{align}
 D&=C(1+A+b+\mu)(d+m+2)(B+2)\ge4,\\
 W&=(1+L_\psi)(H+2)D^{H+2},\\
 K_0&=CQ,\qquad
 K_1=C(m+1)(B+2)W\ge e.
 \label{eq:resources}
\end{align}
The constants in these definitions are independent of $H,d,m,N,B$.

\begin{theorem}[Sharp expected risk at a fixed likelihood]\label{thm:sharp}
Fix the finite learner architecture before sampling. At
$\beta_*=1/(256s)$, every comparator satisfies
\begin{equation}
 \E_{S_n}R(\bar f_{S_n,\beta_*})
 \le R(g)+\frac{5+256sK_0\log(K_1(n+1)^2)}{n+1}.
 \label{eq:sharp}
\end{equation}
The constants are uniform over every finite $N\ge m$. The infimum over deterministic comparator choices may be taken on the right. This is a statement in expectation over the training sample.
\end{theorem}
\begin{proof}
Lemma~\ref{lem:comparison} constructs a sample-independent law whose expected squared uniform error is at most $r=(n+1)^{-2}$ and whose KL is at most $K_0\log(K_1/r)$. Its risk is at most $R(g)+4\sqrt r+r$, because $g$ and $f_0$ lie in $[-1,1]$ and the conditional noise mean is zero. The general final-posterior theorem proved in Section~\ref{sec:sharp-expectation-proof} bounds risk by this comparison risk plus $256s$ times its KL divided by $n+1$. This gives~\eqref{eq:sharp}.
\end{proof}

\begin{theorem}[A high-probability noise-separated oracle]\label{thm:main}
Fix $0<\eta\le1$, $n\ge2$, and use $\beta_\eta=\eta/(32s)$. For any deterministic comparator, with probability at least $1-\delta$,
\begin{align}
 R(\bar f_{S_n,\beta_\eta})
 &\le \int R(f_w)\,d\Pi_{S_n,\beta_\eta}(w)\\
 &\le R_0+(1+\eta)\{R(g)-R_0\}
 +\frac{Cs}{\eta n}\left[K_0\log(K_1n)+\log\frac2\delta\right].
 \label{eq:hp}
\end{align}
The constants are uniform over every finite $N\ge m$. The learner uses neither the comparator nor its complexity parameters.
\end{theorem}

Both likelihood variances are fixed independently of sample size and learner width. They need an upper bound on the noise proxy, and the second also uses the desired approximation slack. The expectation theorem has coefficient exactly one; it is not asserted with high probability.

For fixed activation and output-link constants, the statistical resource is at most
\begin{equation}
 C\big[(d+1)m+Hm^2+1\big]
 \left[(H+2)\log\{C(d+m+2)(B+2)\}+\log\{(H+2)n\}\right].
 \label{eq:explicit-rate}
\end{equation}
Thus both teacher width and depth enter polynomially, and the coefficient cap enters logarithmically. At fixed depth and input dimension, any class with width $o(\sqrt{n/\log n})$ and polynomially bounded coefficients has a vanishing oracle remainder; the single-hidden-layer dependence is only linear in width before logarithms. More generally the exact joint resource~\eqref{eq:explicit-rate} specifies which simultaneously growing widths and depths are allowed.

The high-probability theorem can be made simultaneous over integer comparator resource bounds by a summable confidence allocation. This is only a proof device; no mixture over comparator architectures or sizes is added to the prior. The architecture is fixed before the sample. Uniform constants in $N$ do not assert one common probability event over unrestricted data-dependent choices of infinitely many learner architectures.

\section{A comparison law around a shared network}
Fix a precision $0<\epsilon\le1/4$. All entropy calculations are first made on the latent gates and edge factors. Relative entropy cannot increase when these variables are mapped to actual weights or functions.

\begin{lemma}[Gate and factor costs]\label{lem:scalar-gates}
The following modifications have the stated costs, uniformly in $N$.
\begin{enumerate}
\item A selected gate can be made uniform on $[1-\epsilon,1+\epsilon]$, at KL cost at most $\log N+C\log(C/\epsilon)$.
\item Any specified collection of inactive gates in one layer can be conditioned to have total mass at most $\epsilon$, at joint KL cost at most $C\log(C/\epsilon)$.
\item A Student factor can be made uniform on $[a-\epsilon,a+\epsilon]$, $|a|\le B$, at KL cost at most $C\log(C(B+2)/\epsilon)$.
\end{enumerate}
\end{lemma}
\begin{proof}
Write $a_N=\kappa/N\le1$. On $[3/4,5/4]$, the gamma density
$s^{a_N-1}e^{-s}/\Gamma(a_N)$ is at least $c a_N$, since $\Gamma(a_N+1)$ is bounded above and below on $[0,1]$. This proves the first claim by a pointwise density ratio.

The sum of $k$ inactive gates is $\operatorname{Gamma}(a,1)$ with $a=k\kappa/N\le1$. If $k=0$, no change is needed. Otherwise
\[
 \Pp\{T\le\epsilon\}
 \ge\frac{e^{-\epsilon}\epsilon^a}{\Gamma(a+1)}.
\]
The negative logarithm is at most $C\log(C/\epsilon)$, proving the second claim. In particular, every inactive gate is then at most $\epsilon$.

Finally, the Student density is bounded below by $c(B+2)^{-(\nu+1)}$ on each specified coefficient interval. The uniform interval density gives the last claim.
\end{proof}

For each layer, partition $[N]$ into $m_\ell$ disjoint groups, each of size at least $\lfloor N/m_\ell\rfloor$. Choose one position uniformly from each group. The groups label the teacher neurons. Given these choices, modify the prior as follows:
\begin{itemize}
\item put each selected gate uniformly near one;
\item condition the sum of all other gates in that layer to be at most $\epsilon$;
\item put each selected first-layer row, each factor between two selected neurons, each selected bias, and each selected readout factor uniformly near its teacher coefficient;
\item leave all other edge and bias factors at their native Student laws.
\end{itemize}
There are exactly $P$ modified factors. In particular, a zero coefficient inside an active block must also be controlled. Conditional on a placement, the gate changes are independent of all factor changes, and factor coordinates remain independent.

\begin{lemma}[Entropy independent of learner width]\label{lem:entropy}
The resulting mixture $\rho_\epsilon$ satisfies
\[
 \KL(\rho_\epsilon\|\mu_N)
 \le CQ\log\frac{C(m+1)(B+2)}\epsilon,
 \qquad N\ge m.
\]
\end{lemma}
\begin{proof}
For a fixed placement, Lemma~\ref{lem:scalar-gates} bounds entropy by
\[
 V\log N+CQ\log\frac{C(B+2)}\epsilon.
\]
Distinct placements have disjoint latent supports: a selected gate exceeds $3/4$, while all inactive gates are at most $1/4$. Each group's selected position identifies its teacher role. The KL of the uniform mixture is consequently the average component KL minus the logarithm of the number of placements. For every layer,
\[
 \log N-\log\lfloor N/m_\ell\rfloor\le\log(2m_\ell).
\]
All $V\log N$ costs therefore cancel up to $V\log(2m)$, which is absorbed into the stated bound. This includes $N=m$; no large-width limit is taken. Actual weight vectors need not identify the latent gates. Disjointness is used only before applying the deterministic map to weights.
\end{proof}

\section{Uniform approximation without a robustness restriction}
The comparison suppresses the total inactive gate mass. This gives a direct estimate for its contribution even when individual native hidden features are unbounded.

\begin{lemma}[Uniform second moments and approximation]\label{lem:stability}
For every fixed placement in the comparison construction,
\[
 \left(\E_{\rho_\epsilon}\sup_{x\in\mathcal X}
       |z_w(x)-z_g(x)|^2\right)^{1/2}
 \le\epsilon(H+2)D^{H+2}.
\]
The same statement holds for the mixture over placements.
\end{lemma}
\begin{proof}
Let $R_\ell$ be the maximum, over all learner neurons at layer $\ell$, of the $L^2$ norm of its uniform absolute feature. Let $T_\ell$ be the largest uniform absolute teacher feature, and let $e_\ell$ be the maximum $L^2$ uniform discrepancy at selected neurons.

Conditional on the placement, the gate vector $S^\ell$ is independent of the features $h^\ell$: those features use only gates from earlier layers. The incoming factors at the next layer are also independent of these gates and features. These facts continue to hold when the inactive gates have been jointly conditioned on their total mass.

Write $r_j=\sup_x|h_j^\ell(x)|$. Weighted Cauchy--Schwarz and the inactive-mass bound give, for any next-layer row,
\[
 \sup_x\left|\sum_{j\ \mathrm{inactive}}S_j^\ell G_{ij}h_j^\ell(x)\right|^2
 \le\epsilon\sum_{j\ \mathrm{inactive}}S_j^\ell G_{ij}^2r_j^2.
\]
Taking expectations and using independence bounds this by
$\epsilon^2\mu^2R_\ell^2$. No factor of $N$ remains. The active part has $L^2$ norm at most
$2m\max\{B+1,\mu\}R_\ell$, by Minkowski's inequality. Biases have $L^2$ norm at most $\max\{B+1,\mu\}$.

The first layer has $R_1,T_1\le D^2$. At later layers, the preceding bounds yield
\[
 R_\ell\le b+A\left[
 \max\{B+1,\mu\}+
 \{2m\max\{B+1,\mu\}+\epsilon\mu\}R_{\ell-1}
 \right].
\]
The deterministic teacher has the analogous, smaller bound. The constant in $D$ can therefore be chosen once, independently of depth, so that
\[
 R_\ell,T_\ell\le D^{\ell+1}.
 \label{eq:feature-envelope}
\]

At the first layer, $e_1\le A(d+1)\epsilon$. At an active internal coefficient,
\[
 |S_jG_{ij}-a_{ij}|\le\epsilon(B+2),
 \qquad |S_jG_{ij}|\le2(B+1).
\]
Subtract the teacher row, separate coefficient and feature errors, and include the inactive contribution and the bias perturbation. This gives
\[
 e_\ell\le A\left[
 \epsilon+2m(B+1)e_{\ell-1}
 +m\epsilon(B+2)T_{\ell-1}+\epsilon\mu R_{\ell-1}
 \right].
\]
With the same sufficiently large constant in $D$, induction proves
$e_\ell\le\epsilon(\ell+1)D^{\ell+1}$. The readout has the same estimate without the factor $A$, including its bias. It is bounded by
$\epsilon(H+2)D^{H+2}$, as claimed. Every estimate holds separately for every placement, so averaging preserves it.
\end{proof}

\begin{lemma}[The comparison interface]\label{lem:comparison}
For every $0<r\le1$, there is a sample-independent comparison law $\rho_r$ such that
\[
 \E_{\rho_r}\sup_{x\in\mathcal X}|f_w(x)-g(x)|^2\le r,
 \qquad \KL(\rho_r\|\mu_N)\le K_0\log(K_1/r).
\]
This holds for every finite $N\ge m$, with the resources in~\eqref{eq:resources}.
\end{lemma}
\begin{proof}
Take $\epsilon=\sqrt r/(4W)$. It is at most $1/4$. The Lipschitz property of $\psi$ and Lemma~\ref{lem:stability} prove the approximation bound. Lemma~\ref{lem:entropy} gives the entropy estimate after enlarging the constants in $K_0,K_1$. The construction depends on the comparator and the desired accuracy, but not on any observations.
\end{proof}

The uniform stability factor can grow exponentially with depth. Its logarithm enters the entropy, so this does not create exponential sample complexity. The argument requires no stability assumption on the teacher or on the support of the prior.

\section{A nearly linear depth cost for contractive comparison rows}
The unconditional theorem permits arbitrary bounded coefficients. A natural condition on the comparison network alone improves its depth cost further. The prior remains exactly the same and is not conditioned on this property.

\begin{theorem}[Contractive internal teacher rows]\label{thm:stable-teachers}
Suppose every internal hidden row of the comparator satisfies
\[
 A\sum_j|a^\ell_{ij}|\le1\qquad(2\le\ell\le H).
\]
The first layer, all biases, and the readout may have arbitrary coefficients bounded by $B$. Then all comparison, learning and uniform-neighborhood conclusions hold with $W$ replaced by
\[
 W_{\rm st}=C(1+L_\psi)(d+1)(m+1)^2(B+1)^2(H+1)^2,
\]
and $K_1$ replaced by $C(m+1)(B+2)W_{\rm st}$. Here $C$ may depend on the fixed activation and scalar laws. In particular the statistical resource is
\[
 C\big[(d+1)m+Hm^2+1\big]
       \log\{C(d+1)(m+1)(B+1)(H+1)n\}.
\]
For fixed input dimension this has dominant width--depth dependence
$\widetilde O(m^2H)$.
\end{theorem}
\begin{proof}
Use exactly the same comparison laws as before, with precision
$\epsilon\le c/[(m+1)(H+1)]$. Let $S_\ell$ denote the largest selected feature's $L^2$ uniform magnitude, and let $P_\ell$ denote the largest such magnitude over all learner features. Put $R=C(d+1)(B+1)$, with constants allowed to depend on the fixed activation and scalar laws. The first layer has $S_1,P_1\le R$.

For a selected row, write its effective coefficient as
$S_jG_{ij}=a_{ij}+\Delta_{ij}$. Since $|S_j-1|\le\epsilon$,
$|G_{ij}-a_{ij}|\le\epsilon$, and $\sum_j|a_{ij}|\le1/A$,
\[
 A\sum_j|\Delta_{ij}|\le C(m+1)\epsilon.
\]
The inactive contribution has $L^2$ uniform norm at most $\mu\epsilon P_{\ell-1}$, as before. Biases add at most a fixed multiple of $R$. Thus
\begin{align*}
 S_\ell&\le CR+\{1+C(m+1)\epsilon\}S_{\ell-1}
                  +C\epsilon P_{\ell-1},\\
 P_\ell&\le CR+C(m+1)S_{\ell-1}+C\epsilon P_{\ell-1}.
\end{align*}
The second inequality bounds native rows using their finite factor moment and includes selected rows by enlarging $C$.

If $C\epsilon\le1/2$, the second recurrence gives
$P_\ell\le C(m+1)\max\{R,S_1,\ldots,S_\ell\}$. Substitution in the first and discrete Gronwall give
\[
 \max_{\ell\le H}S_\ell\le CR(H+1),\qquad
 \max_{\ell\le H}P_\ell\le C(m+1)R(H+1),
\]
provided the constant in the precision restriction is sufficiently small.

For discrepancies, keep the deterministic teacher coefficients on the previous error and put the coefficient perturbations on the actual previous features. Contractivity then gives
\[
 e_\ell\le e_{\ell-1}
       +C\epsilon\{1+(m+1)S_{\ell-1}+P_{\ell-1}\}.
\]
Starting from $e_1\le A(d+1)\epsilon$ and summing yields
$e_H\le C\epsilon(m+1)R(H+1)^2$.
The readout's deterministic coefficient sum is at most $mB$. Its own gate/factor errors and inactive contribution therefore give
\[
 \left(\E\|z-z_g\|_\infty^2\right)^{1/2}
 \le C\epsilon(d+1)(m+1)^2(B+1)^2(H+1)^2.
\]
Multiplication by the output-link Lipschitz constant proves the claimed stability factor. Taking $\epsilon=\sqrt r/(4W_{\rm st})$, with a sufficiently large fixed constant in $W_{\rm st}$, also enforces the precision restriction. The entropy proof is unchanged, and all statistical and prior-mass consequences follow by substitution.
\end{proof}

This condition is a statement about a useful implementation of the comparator. It is not a claim that every small Boolean circuit has such an implementation, or that the Bayesian posterior is supported on contractive networks.

\section{A sharper resource for Boolean circuits}
\label{sec:ard-Boolean-restoration}
Boolean circuits admit comparison implementations whose errors are restored at each logical gate. Under the neuron-gated prior this removes the extra depth factor in the general sensitivity bound, without restricting learner weights.

Consider a layered fan-in-two Boolean circuit, counting all COPY and delay gates explicitly. Let $S$ be its gate count, $h$ its depth and $C$ its maximum layer width. Its input bits are the signs of $x\in\{-1,1\}^d$. There are explicit ReLU, tanh and sigmoid implementations with a bounded final clipping link and numerical coefficient caps. ReLU needs hidden depth $2h$ and width at most $2C$; tanh and sigmoid need depth $h$ and width at most $C$. The implementations pad to larger depths $H$. Let $m$ be their actual maximum width, let $V$ count their hidden vertices, and let $P$ count their dense active blocks and biases as in the general comparison theorem. Put $Q=P+V+H$. In particular,
\[
 Q\le C_0\{(d+1)m+m^2H+H\},
\]
with a numerical constant $C_0$.

\begin{theorem}[Boolean comparison at nearly parameter-count cost]
\label{thm:ard-Boolean-restoration}
For every such implementation and every learner width $N\ge m$, the neuron-gated prior has a data-independent comparison law with
\[
 \E_\rho\sup_{x\in\{-1,1\}^d}|f_w(x)-f(x)|\le t,
 \qquad
 \KL(\rho\|\mu_N)\le C Q\log\frac{C Q(m+1)}t
 \quad(0<t\le1).
\]
It also assigns the correct raw sign with a fixed positive all-input margin prior mass at least $\exp[-CQ\log(CQ(m+1))]$. Consequently the sharp expected-risk remainder for comparison with this Boolean function is
\[
 \frac{C(1+\sigma^2)Q\log\{C(n+1)Q(m+1)\}}{n+1},
\]
independently of learner width. The exact Bayesian sign learner has total confidence and majority mistake bounds of the same nearly linear-$Q$ order, along arbitrary input sequences.
\end{theorem}
\begin{proof}
For ReLU encode bits by zero and one. Use two hinges followed by their ReLU difference:
\[
 \begin{array}{c|c|c}
 &r&s\\ \hline
 \mathrm{AND}(a,b)&\operatorname{ReLU}(4(a+b)-5)&
                    \operatorname{ReLU}(4(a+b)-6)\\
 \mathrm{NOT}(a)&\operatorname{ReLU}(2-4a)&
                    \operatorname{ReLU}(1-4a)\\
 \mathrm{COPY}(a)&\operatorname{ReLU}(4a-1)&
                    \operatorname{ReLU}(4a-2).
 \end{array}
\]
The next hidden layer computes $\operatorname{ReLU}(r-s)$. In the first logical layer substitute $(1+x_i)/2$ directly into these affine expressions, requiring no input-feature layer. The final readout is $2a-1$. Extra hidden depth may relay the final nonnegative bit.

There are numerical constants $u_0,K_0$ such that a total uniform preactivation perturbation of size at most $u\le u_0$ over all selected rows and the readout keeps the raw output within $C u$ of the target, and all selected hidden values below $K_0$. Indeed the hinge margins preserve which pair is active. Both-positive hinge differences reset the new bit to one plus that gate's perturbations; both-negative pairs reset it to zero. No previous bit error is amplified once it lies inside the fixed decision margin. Extra identity relays add only their portion of the total perturbation budget.

Fix an embedding in the general comparison law at scalar precision $\epsilon$. For analysis only, cap selected ReLU activations to $[0,K_0+1]$. Let $R_\ell$ be the largest expected uniform magnitude of a modified unselected activation. Its active-parent contribution is at most $C m$, since selected gates are bounded by two, selected features are capped, and each unmodified edge factor has finite first absolute moment. The inactive gate sum is at most $\epsilon$, independently of its own incoming features. Native biases have finite first moment. Hence
\[
 R_\ell\le C(m+1)+C\epsilon R_{\ell-1}.
\]
The first layer has bound $C(d+1)$, so at the first hidden transition one may retain $C\epsilon(d+1)$; thereafter $R_\ell\le C(m+d+1)$ whenever $C\epsilon\le1/2$. This harmless input term will be absorbed by $Q$ below.

At a selected row, each active product coefficient differs from the teacher coefficient by at most $C\epsilon$, including coefficients that are zero in the teacher. The $m$ capped active features therefore contribute at most $C m\epsilon$ to its uniform residual. The selected bias contributes $\epsilon$. The inactive contribution has expected uniform magnitude at most $C\epsilon R_{\ell-1}$. First-layer residuals have size at most $(d+1)\epsilon$. Thus the expected sum of all selected-row and readout residual magnitudes is at most
\[
 C(V+1)(m+d+1)\epsilon\le C Q^2\epsilon.
\]
On the event that this sum is at most $u\le u_0$, restoration proves inductively that the cap is never reached, so the modified and actual computations coincide. Their raw output error is at most $Cu$. Since the final clipping link bounds all outputs, Markov's inequality gives expected uniform clipped error at most
$Cu+C Q^2\epsilon/u$. Take $u=ct$ and $\epsilon=c t^2/Q^2$, reducing the constants if necessary. The general entropy bound is then at most $C Q\log(CQ(m+1)/t)$. A fixed $u$ and $\epsilon=c/Q^2$ give the raw margin probability, by the same event and binary relative entropy.

For tanh, encode signs in $[0.9,1]$ and $[-1,-0.9]$, and use preactivations $4(a+b-1)$, $-4a$ and $4a$ for AND, NOT and COPY. For sigmoid, encode bits in $[0,0.1]$ and $[0.9,1]$, and use $20(a+b-1.5)$, $20(0.5-a)$ and $20(a-0.5)$. Substitute the exact input-bit encoding into first-layer expressions. These gates preserve their encoding intervals under a preactivation error of magnitude $0.1$. The final readouts $2a$ for tanh and $4a-2$ for sigmoid clip exactly to the target signs on this event. Every native feature is bounded, so each selected-row residual has expected uniform magnitude at most $C(m+d+1)\epsilon$ directly; a union bound suffices. The same entropy conclusion follows, with no auxiliary cap.

The expected-risk conclusion follows by applying the sharp all-data theorem at $t=(n+1)^{-1}$. The margin-mass conclusion and exact-sign online argument give the remaining assertions, with the usual nonanticipating requirement for posterior-sampled mistakes.
\end{proof}

\section{Nearly additive extra learner depth for ReLU}
A shallow comparison function can also be learned by a much deeper architecture. With biases available, clipping can be computed through a single scalar channel:
\begin{equation}
 \operatorname{clip}(z)=1-\operatorname{ReLU}\!\left(2-\operatorname{ReLU}(z+1)\right).
 \label{eq:scalar-reset}
\end{equation}
This identity permits a bounded scalar relay after the comparison computation, without increasing its maximum width.

\begin{theorem}[A scalar relay through extra layers]\label{thm:scalar-padding}
Use ReLU hidden activations and final clipping in the neuron-gated prior. Let a comparison network have hidden depth $h\ge1$, maximum width $m$, coefficient cap $B\ge1$, and resource
\[
 Q_g=P_g+V_g+h+1,
\]
where $P_g,V_g$ are the counts in~\eqref{eq:counts} for its own depth. Prescribe any learner depth $H\ge h$ and width $N\ge m$. Put
\[
 D_0=C(d+m+2)(B+4),\qquad L=(h+4)\log D_0+\log(h+4),
\]
and
\begin{equation}
 J_n=C\left[
 Q_g\{L+\log((m+1)(n+1))\}
 +(H-h+1)\log\{C(H+1)(n+1)\}\right].
 \label{eq:scalar-padding-cost}
\end{equation}
At the same fixed likelihood temperature $\beta_*=1/(256s)$,
\[
 \E R(\bar f_{S_n,\beta_*})
 \le R(g)+\frac{C s(1+J_n)}{n+1}.
\]
The high-probability noise-separated companion holds with remainder
$Cs[J_n+\log(2/\delta)]/(\eta n)$ at $\beta_\eta=\eta/(32s)$.
The extra learner-depth cost is nearly linear and additive; it is not multiplied by $m^2$ or by the logarithm of the teacher coefficient cap. No extra width allowance or prior restriction is needed.
\end{theorem}
\begin{proof}
For $H=h$, use the general comparison directly. If $H=h+1$, one added neuron computes $\operatorname{ReLU}(z_g+1)$, with readout that subtracts one; final clipping still gives $g$. The general comparison at this depth has the stated resource. Hence suppose $H\ge h+2$ and put $q=H-h-2$.

Append the two scalar ReLU neurons in~\eqref{eq:scalar-reset}. The second has ideal feature $u(x)\in[0,2]$, and the final intended raw output is $1-u(x)=g(x)$. Insert $q$ identity ReLU relays between this feature and its readout. All added layers have width one. The original readout coefficients become the first reset neuron's incoming coefficients; the second reset has coefficient $-1$ and bias $2$. Every subsequent relay has coefficient $1$ and bias zero. This implementation uses at most $P_g+2(q+2)$ coefficient positions and $V_g+q+2$ hidden vertices.

Fix a desired expected uniform error $0<t\le1$. In the teacher and reset block use the comparison construction at a common precision
\[
 \epsilon_0=ct\big/\{(h+4)D_0^{h+4}\},
\]
with a sufficiently small fixed $c$. The outgoing gates of the final reset neuron are not yet specified; they do not affect its own feature. The uniform stability proof gives
\[
 e_0:=\E\|h^0-u\|_\infty\le t/100,
 \qquad T_0:=\E\|h^0\|_\infty\le2+t/100,
 \qquad M_0\le U,
\]
where $M_0$ is the largest expected uniform magnitude of any native or selected feature at that layer, and $\log(U+1)\le CL$. The entropy paid by this prefix is at most
$CQ_g[L+\log(C(m+1)/t)]$.

For each relay, use a precision $\epsilon_\ell$ for its selected incoming factor and bias, and for the gates in the preceding layer. Thus its selected gate is uniform near one, and the total inactive gate mass is at most $\epsilon_\ell$. The layerwise choices remain independent of the features whose outgoing gates they modify. The designated coefficient and gate are positive with means one, so their product has expected absolute value exactly one. If $T_\ell$ is the selected feature's expected uniform magnitude and $M_\ell$ is the maximum over all features, then
\begin{align}
 T_\ell&\le T_{\ell-1}+C\epsilon_\ell(1+M_{\ell-1}),\\
 M_\ell&\le C(1+T_{\ell-1})+C\epsilon_\ell M_{\ell-1}.
 \label{eq:scalar-relay-moments}
\end{align}
The first inequality uses the selected bias and inactive background. The second uses the native factor first moment; it also covers the selected row by enlarging $C$.

For $q\ge1$, choose
\[
 \epsilon_1=\frac{ct}{(q+1)(1+U)},\qquad
 \epsilon_\ell=\frac{ct}{q+1}\quad(2\le\ell\le q),
\]
with $c$ small. The first relay makes $M_1$ bounded by a numerical constant. Induction in~\eqref{eq:scalar-relay-moments}, summing the first recurrence and using $C\epsilon_\ell\le1/2$, gives $T_\ell\le3$ and $M_\ell\le C$ for all later relays. The identity target and independence of each incoming gate/factor from the preceding feature similarly give
\[
 e_\ell\le e_{\ell-1}+C\epsilon_\ell(1+M_{\ell-1}),
\]
so their total error contribution is at most $Ct$ with an arbitrarily small fixed coefficient after choosing $c$.

At the final readout, use target coefficient $-1$ and bias $1$. If $q\ge1$, its precision and the outgoing gate precision are $ct/(q+1)$. If $q=0$, use $ct/(1+U)$ instead. The same first-moment calculation bounds its added expected uniform raw error by a small multiple of $t$. The whole comparison consequently has expected uniform raw error at most $t$ from $g$.

The entropy argument permits distinct deterministic precisions at different layers. The selected gate markers still identify every chosen vertex, and its location entropy cancels its $\log N$ cost. Each scalar relay layer has a bounded number of modified factors and gates. Their combined entropy, including the readout, is at most
\[
 C(q+1)\log\frac{C(q+1)}t+C\log(1+U).
\]
Only the first relay, or the readout when there is no relay, pays for the large prefix background. Adding the prefix cost gives~\eqref{eq:scalar-padding-cost} with $n+1$ replaced by $1/t$.

Final clipping cannot enlarge the approximation error. Bounded predictions and the bounded regression mean imply that comparison risk exceeds $R(g)$ by at most $4t$. Take $t=(n+1)^{-1}$ and apply the sharp expectation theorem. For the high-probability version, bounded differences also give
$\E_\rho\|f-g\|_{L^2}^2\le2t$; the relative-loss proof with $t=1/n$ then gives the stated companion.
\end{proof}

For fixed teacher depth and coefficient bound, this gives a resource of order
$\widetilde O(d m+m^2h+H)$, rather than multiplying the large teacher width by every added learner layer. The likelihood and prior do not know the comparator depth or use a mixture over depths.

\begin{theorem}[A scalar relay for tanh]\label{thm:tanh-padding}
Use tanh hidden activations, any fixed bounded Lipschitz output link, and the Student factor law. Let the teacher have depth $h$, width $m$, coefficient bound $B$, and resource $Q_g$ as above. For every learner depth $H\ge h$ and every finite width $N\ge m$, the sharp expected and near-sharp high-probability conclusions of Theorem~\ref{thm:scalar-padding} hold with
\begin{align}
 J_n^{\tanh}=C\big[&Q_g\{(h+2)\log D_0+
                       \log((H+1)(n+1))\}\notag\\
 &+(H-h+1)\log\{C(H+1)(m+1)(B+2)(n+1)\}\big],
 \label{eq:tanh-padding-cost}
\end{align}
where $D_0=C(d+m+2)(B+2)$. Extra learner depth again has only a nearly linear additive cost. The fast rate in this refinement uses the Student factor's logarithmic cost for a large final readout; it is not asserted for the fully Gaussian-factor variant.
\end{theorem}
\begin{proof}
The case $H=h$ is the general theorem. Put $q=H-h\ge1$ and $M=(m+1)B$, so the teacher raw output satisfies $|z_g|\le M$ everywhere. For a desired uniform error $t$, choose
\[
 \delta=\sqrt{ct/(qM^3)},
\]
with a small fixed constant, absorbing the output-link Lipschitz constant in $c$. Append $q$ scalar tanh layers: the first computes $u_1=\tanh(\delta z_g)$, and the others compute $u_\ell=\tanh(u_{\ell-1})$. Their intended final readout is $u_q/\delta$. Since $|u_\ell|\le\delta M$ and $|\tanh v-v|\le |v|^3/3$,
\[
 \|u_q/\delta-z_g\|_\infty\le q\delta^2M^3/3\le ct.
\]
The first new row has coefficients $\delta v_j$ and bias $\delta b$, all bounded by one. The remaining rows have coefficient one and bias zero. Width remains at most $m$.

Use the original comparison construction in the teacher block with precision small enough that the largest selected feature error is at most $t/(100M)$. Its entropy is at most
$CQ_g[(h+2)\log D_0+\log(C/t)]$. As before, the teacher's final outgoing gates can be specified later, since they do not enter its own features.

For the first scalar row and all subsequent scalar hidden rows, use factor and preceding-gate precision
\[
 \epsilon_s=\frac{ct\delta}{q(m+1)}.
\]
All actual tanh features are bounded by one. The first scalar preactivation discrepancy is therefore at most
$\delta M$ times the teacher feature error plus $C(m+1)\epsilon_s$ in expected uniform norm. At later unit relays, the selected positive gate and coefficient have product with expected absolute value one, independently of the preceding feature. If $e_\ell$ denotes the expected uniform discrepancy from $u_\ell$, then
\[
 e_\ell\le e_{\ell-1}+C\epsilon_s.
\]
The selected biases and the inactive gate mass are included in this additive term. Consequently $e_q\le c_1t\delta$, with $c_1$ arbitrarily small after reducing the comparison precision constants.

At the final readout use coefficient $1/\delta$, bias zero, and factor and gate precision $\epsilon_o=ct/(M+1)$. The selected coefficient is positive, and the gate has mean one, so propagation of $e_q$ costs exactly $e_q/\delta$ in expected absolute norm. Its coefficient perturbation acting on $|u_q|\le\delta M$ contributes at most $C\epsilon_oM$. The bias and inactive terms contribute at most $C\epsilon_o$, since tanh features are bounded. Together with the deterministic relay error this gives expected uniform raw error at most $t$.

The entropy cancellation still applies layer by layer. The first scalar row and the teacher's outgoing gates use $O(m)$ controlled variables; the other relays use $O(q)$; and the final readout coefficient contributes only $C\log(1/\delta)$. Since
\[
 \log(1/\delta)\le C\log\{CqM/t\},
\]
the additional entropy is at most
$C(m+q+1)\log\{Cq(m+1)M/t\}$.
Absorb the $m$ term into $Q_g$ and combine with the prefix cost to obtain~\eqref{eq:tanh-padding-cost}. The same statistical substitutions as in the ReLU proof finish both conclusions.
\end{proof}

\section{The fixed-temperature noisy oracle}
Write $A(f)=\E(f(X)-f_0(X))^2=R(f)-R_0$, $s=1+\sigma^2$ and $\beta=\eta/(32s)$. Network outputs and the regression mean lie in $[-1,1]$, but observed labels need not be bounded. The following two elementary inequalities provide the statistical part of the proof.

\begin{lemma}[A relative-loss PAC--Bayes inequality]\label{lem:relative-pac}
With probability at least $1-\delta/2$, simultaneously for every probability measure $\rho$ on weights,
\begin{equation}
 c_\beta^{\mathrm{rel}}\E_\rho A(f_w)
 \le\E_\rho[\whR_S(f_w)-\whR_S(f_0)]
       +\frac{\KL(\rho\|\mu_N)+\log(2/\delta)}{\beta n},
 \label{eq:relative-pac}
\end{equation}
where
\[
 c_\beta^{\mathrm{rel}}=\frac{1-e^{-4\beta(1-2\beta\sigma^2)}}{4\beta}
       \ge1-4s\beta=1-\eta/8.
\]
\end{lemma}
\begin{proof}
Fix $f=f_w$ and put $e=f-f_0$. The difference of squared losses is $e(X)^2-2\xi e(X)$. Conditional sub-Gaussianity gives
\[
 \E[e^{-\beta\{(f(X)-Y)^2-(f_0(X)-Y)^2\}}\mid X]
 \le e^{-\beta(1-2\beta\sigma^2)e(X)^2}.
\]
Since $0\le e^2\le4$, convexity on that interval yields
\[
 \E e^{-\beta\{(f(X)-Y)^2-(f_0(X)-Y)^2\}}
 \le e^{-\beta c_\beta^{\mathrm{rel}} A(f)}.
\]
The lower bound on $c_\beta^{\mathrm{rel}}$ follows from $1-e^{-u}\ge u-u^2/2$ and the chosen range of $\beta$. Independence, integration against the prior, Markov's inequality and the entropy variational inequality prove~\eqref{eq:relative-pac}. The individual squared losses have finite expectation under the conditional sub-Gaussian noise assumption; their difference has the exponential moment just established.
\end{proof}

The comparison measure may itself be averaged before applying concentration. This is useful because the perturbation lemma is naturally a statement about its expected risk.

\begin{lemma}[Concentration for a fixed comparison law]\label{lem:comparison-concentration}
Let $\nu$ be any fixed probability measure on bounded predictions, independent of the data, and let $a=\E_\nu A(f_w)$. With probability at least $1-e^{-\ell}$,
\begin{equation}
 \E_\nu[\whR_S(f_w)-\whR_S(f_0)]
 \le a+C\sqrt{sa\ell/n}+Cs\ell/n.
 \label{eq:comparison-bernstein}
\end{equation}
\end{lemma}
\begin{proof}
Set
\[
 v(x)=\E_\nu(f_w(x)-f_0(x))^2,\qquad
 b(x)=\E_\nu(f_w(x)-f_0(x)).
\]
Then $0\le v\le4$, $b^2\le v$, and one comparison-averaged excess loss is $L=v(X)-2\xi b(X)$, of mean $a$. For $0\le t\le1/(16s)$, conditional sub-Gaussianity gives
\[
 \E e^{t(L-a)}
 \le e^{-ta}\E e^{(t+2\sigma^2t^2)v(X)}
 \le\exp\left\{a\left[\frac{e^{4(t+2\sigma^2t^2)}-1}{4}-t\right]\right\}
 \le e^{6sa t^2}.
\]
For the last inequality, put $u=t+2\sigma^2t^2\le9t/8$, use $4u\le1$ and $e^{4u}-1\le4u+16u^2$. Chernoff's bound for the independent observations then gives~\eqref{eq:comparison-bernstein}, with an absolute constant. Equivalently, optimize over $t\le1/(16s)$ to obtain a square-root term and a linear term.
\end{proof}

\begin{proof}[Proof of Theorem~\ref{thm:main}]
Fix the deterministic comparator $g$. For $0<r\le1$, let $Q$ be the comparison law of Lemma~\ref{lem:comparison}. Weighted triangle gives
\[
 a:=\E_Q A(f_w)
 \le(1+\eta/8)A(g)+(1+8/\eta)r.
\]
At the Gibbs posterior, the empirical-risk plus relative-entropy term equals the free energy
$F_{S_n,\beta}=-(\beta n)^{-1}\log Z_{S_n,\beta}$.
The variational formula and Lemma~\ref{lem:relative-pac} therefore yield
\[
 c_\beta^{\rm rel}\E_{\Pi_{S_n,\beta}}A(f_w)
 \le \E_Q[\whR_{S_n}(f_w)-\whR_{S_n}(f_0)]
       +\frac{K_0\log(K_1/r)+\log(2/\delta)}{\beta n}.
\]
Use Lemma~\ref{lem:comparison-concentration} with $\ell=\log(2/\delta)$.
Young's inequality bounds its square-root term by
$(\eta/8)a+Cs\ell/(\eta n)$.
The coefficient of $A(g)$ after division by
$c_\beta^{\rm rel}\ge1-\eta/8$ is at most
\[
 \frac{(1+\eta/8)^2}{1-\eta/8}\le1+\eta.
\]
Every other term is at most
\[
 \frac C\eta\left[
 r+\frac{s}{n}\{K_0\log(K_1/r)+\log(2/\delta)\}\right].
\]
Take $r=1/n$. Since $K_0\ge1$, $K_1\ge e$ and $s\ge1$, the extra $1/n$ is absorbed. Add $R_0$ and apply Jensen's inequality to the posterior mean. The two concentration events fail with total probability at most $\delta$.
\end{proof}

The regression function $f_0$ and its empirical risk appear only in the proof; their terms cancel from the learner's variational comparison. The actual posterior uses only the observations, the chosen approximation factor and an upper bound on the conditional noise proxy.

% Independent proof module, 18 November 2026.
% This file is a review artifact, not an edit to an owner manuscript.
% The leave-one-out strategy is motivated by Hogsgaard--Rebeschini--Wegel,
% arXiv:2607.02247 (2026). Their stated theorem is for finite dictionaries
% and bounded losses. The arbitrary-prior and conditional sub-Gaussian
% extension below is derived here; it should receive independent checking.

\section{A sharp expectation oracle for the ordinary final posterior}
\label{sec:sharp-expectation-proof}

Let $\pi$ be any probability prior on jointly measurable functions
$f_w:\mathcal X\to[-1,1]$. Suppose the observations are IID,
$f_0(x)=\mathbb E[Y\mid X=x]\in[-1,1]$, and
$\xi=Y-f_0(X)$ is conditionally sub-Gaussian with proxy $\sigma^2$.
Set $s=1+\sigma^2$ and fix $T\ge256s$.
For a sample $S_n$ of size $n$, let
\[
 dQ_{S_n}(w)=Z_{S_n}^{-1}
   \exp\left\{-\frac1T\sum_{i=1}^n(f_w(X_i)-Y_i)^2\right\}d\pi(w),
 \qquad \bar f_{S_n}=\int f_w\,dQ_{S_n}(w).
\]
This is the ordinary posterior with Gaussian observation variance $T/2$,
using all $n$ observations. No averaging over sample sizes is performed.

\begin{theorem}[Sharp all-data expectation bound]\label{thm:generic-expectation}
For every deterministic comparison law $\rho\ll\pi$,
\[
 \mathbb E_{S_n} R(\bar f_{S_n})
 \le \int R(f_w)\,d\rho(w)
        +\frac{T\operatorname{KL}(\rho\|\pi)}{n+1}.
\]
Consequently the infimum over all such comparison laws may be taken on
the right-hand side. The theorem requires neither convexity of the
support of $\pi$ nor a Bernstein condition relative to the comparator.
\end{theorem}

\begin{proof}
We first prove an averaged one-observation inequality. Fix a probability
law $q$ on $U\in[-1,1]$, with mean $\mu$ and variance $v$. For a real
label $y$, let $L_y(u)=(u-y)^2$ and
\[
 dp_y(u)=\frac{e^{-L_y(u)/T}}{\int e^{-L_y/T}\,dq}\,dq(u).
\]
For two independent draws $U,U'$ from $q$, the covariance identity gives
\begin{align*}
 \int L_y\,dq-\int L_y\,dp_y
 &=\frac{\mathbb E[(L_y(U)-L_y(U'))
         (e^{-L_y(U')/T}-e^{-L_y(U)/T})]}{2\int e^{-L_y/T}\,dq}.
\end{align*}
Since
\[
 |L_y(u)-L_y(u')|\le2(1+|y|)|u-u'|,
 \qquad \operatorname{osc}_{[-1,1]}L_y\le4(1+|y|),
\]
the mean-value theorem bounds the last display by
\[
 \frac{4(1+|y|)^2}{T}
       e^{4(1+|y|)/T}v.
\]
Also $\int L_y\,dq=L_y(\mu)+v$. Hence
\begin{equation}
 \int L_y\,dp_y\ge L_y(\mu)+
 \left[1-\frac{4(1+|y|)^2}{T}e^{4(1+|y|)/T}\right]v.
 \label{eq:one-step-pointwise}
\end{equation}

Now let $Y=a+\xi$, where $|a|\le1$ and $\xi$ is centered
$\sigma^2$-sub-Gaussian, independently of the fixed law $q$.
The elementary estimates
\[
 \mathbb E e^{t|\xi|}\le2e^{\sigma^2t^2/2},
 \qquad \mathbb E\xi^4\le16\sigma^4
\]
and Cauchy--Schwarz imply
\begin{align*}
 &\mathbb E\left[(1+|Y|)^2e^{4(1+|Y|)/T}\right]\\
 &\quad\le e^{8/T}
       \mathbb E[(8+2\xi^2)e^{4|\xi|/T}]\\
 &\quad\le e^{8/T}
       \left[16e^{8\sigma^2/T^2}
              +8\sqrt2\,\sigma^2e^{16\sigma^2/T^2}\right]
 \le32s
 \qquad(T\ge256s).
\end{align*}
Thus the expected coefficient subtracted in
\eqref{eq:one-step-pointwise} is at most $128s/T\le1/2$.
After averaging over $Y$ we obtain
\begin{equation}
 \mathbb E_Y L_Y(\mu)
 \le\mathbb E_Y\int L_Y\,dp_Y-\tfrac12v
 \le\mathbb E_Y\int L_Y\,dp_Y.
 \label{eq:averaged-one-step}
\end{equation}
The same statement holds conditionally on any input $X=x$ and on an
independent training sample.

Introduce $n+1$ IID observations, retaining the same fixed architecture,
prior and $T$ used by the $n$-observation learner. Let $Q^+$ be their full
posterior and $Q^{-i}$ the posterior after removing observation $i$.
Conditional on the other observations and $X_i$, push $Q^{-i}$ forward
through $w\mapsto f_w(X_i)$ to obtain $q$. It is independent of $Y_i$
given this conditioning. Updating $q$ by the omitted squared-loss factor
gives exactly the corresponding pushforward of $Q^+$.
Applying \eqref{eq:averaged-one-step} and then averaging yields
\[
 \mathbb E (\bar f^{-i}(X_i)-Y_i)^2
 \le\mathbb E\int(f_w(X_i)-Y_i)^2\,dQ^+(w).
\]
Sum over $i=1,\ldots,n+1$. By exchangeability and independence, the
left-hand average is $\mathbb E_{S_n}R(\bar f_{S_n})$.
For each realized enlarged sample the Gibbs variational inequality gives
\begin{align*}
 &\int\sum_{i=1}^{n+1}(f_w(X_i)-Y_i)^2\,dQ^+(w)\\
 &\quad\le\int\sum_{i=1}^{n+1}(f_w(X_i)-Y_i)^2\,d\rho(w)
             +T\operatorname{KL}(\rho\|\pi),
\end{align*}
where the nonnegative term $T\operatorname{KL}(Q^+\|\pi)$ was dropped.
Take expectations and divide by $n+1$ to prove the theorem.
All integrals are finite: the outputs are bounded and the labels have
finite second moment. The posterior normalizers are positive, and the
loss oscillation over parameters is finite for every realized sample.
\end{proof}

\begin{corollary}[Logarithmic comparison entropy]
Suppose a deterministic $g:\mathcal X\to[-1,1]$ has, for every $0<r\le1$,
a data-independent comparison law with
\[
 \mathbb E_\rho\|f_w-g\|_{L^2(P_X)}^2\le r,
 \qquad \operatorname{KL}(\rho\|\pi)\le K_0\log(K_1/r),
 \qquad K_0\ge1,\ K_1\ge e.
\]
At $T=256s$ the same all-data final posterior mean satisfies
\[
 \mathbb E R(\bar f_{S_n})
 \le R(g)+\frac{5+256sK_0\log(K_1(n+1)^2)}{n+1}.
\]
Thus the oracle is sharp, with coefficient one on comparator risk and
a remainder of order $sK_0\log(K_1n)/n$.
\end{corollary}
\begin{proof}
Write $A(g)=\mathbb E(g-f_0)^2\le4$. Cauchy--Schwarz gives
\[
 \int R(f_w)\,d\rho
 \le R(g)+2\sqrt{A(g)r}+r.
\]
Take $r=(n+1)^{-2}$ and use the theorem. The approximation increase is
at most $4/(n+1)+1/(n+1)^2\le5/(n+1)$.
\end{proof}

\paragraph{Scope.}
The result is in expectation over the training sample. It does not
assert a sharp high-probability oracle, and an all-data refit is not
being inferred from a holdout result. Every observation enters the
likelihood of the returned posterior. If the comparison constants are
uniform in learner width, so is this bound. One may take the infimum
over all deterministic comparator classes and resources directly;
unlike a common high-probability event, this expectation oracle needs
no confidence allocation over classes. The learner must still have
its architecture and prior fixed before observing the data.

\paragraph{Attribution.}
The leave-one-out strategy and a pointwise inverse-tilting inequality
for bounded losses are from M.~M.~H\o gsgaard, P.~Rebeschini and
T.~Wegel, \emph{Aggregation with Exponential Weights is Optimal in
Expectation}, arXiv:2607.02247, July 2026~\cite{hogsgaard}. Their main statement uses a
finite uniform dictionary and bounded losses. Above, sums are replaced
by integrals and the variational comparison is written in KL form;
the pointwise bounded-loss step is replaced by an averaged cooling
inequality for bounded predictions and conditionally sub-Gaussian
labels. The arbitrary-prior and conditional sub-Gaussian extension is proved here.

\section{Uniform function mass and online sign prediction}
The comparison law also yields a lower bound on the prior probability of a uniform function neighborhood.

\begin{theorem}[Uniform function-neighborhood mass]\label{thm:uniform-mass}
For every comparator and every $0<t\le1$,
\[
 -\log\mu_N\left\{\sup_{x\in\mathcal X}|f_w(x)-g(x)|\le t\right\}
 \le C Q\log\frac{C(m+1)(B+2)W}{t}.
\]
The same statement holds for raw outputs, using
$W_{\rm raw}=(H+2)D^{H+2}$ in place of $W$.
\end{theorem}
\begin{proof}
Use the comparison with $\epsilon=t/(2W)$. Its expected squared uniform error is at most $t^2/4$, so the neighborhood event has comparison probability at least $3/4$. For any event $E$,
\[
 \KL(\rho\|\mu_N)\ge\rho(E)\log\frac1{\mu_N(E)}-\log2.
\]
Apply this inequality and Lemma~\ref{lem:entropy}. The raw-output argument omits the output-link factor.
\end{proof}

Let a raw teacher have a global sign margin on a prescribed compact input set:
\[
 y(x)z_g(x)\ge2a\quad(x\in\mathcal X),
 \qquad y:\mathcal X\to\{-1,+1\},\quad a>0.
\]
The raw neighborhood with radius $t=\min\{1,a/2\}$ is contained in the set $E_y$ of networks with $y(x)z_w(x)\ge a$ everywhere. Thus $\mu_N(E_y)\ge e^{-J}$, where
\[
 J=CQ\left[(H+2)\log\{C(d+m+2)(B+2)\}
           +\log\{(H+2)(1+1/a)\}\right].
\]
Initially conditioning the prior on the all-input absolute-margin event can only increase this lower bound.

\begin{corollary}[Total online mistakes]\label{cor:online}
Use the preceding prior, optionally conditioned on its all-input margin event, and update by exact agreement with each revealed sign. If $p_t$ is the posterior probability of the correct label before round $t$, then every input sequence satisfies
\[
 \sum_t-\log p_t\le J,\qquad \sum_t(1-p_t)\le J.
\]
Majority prediction therefore makes at most $J/\log2$ mistakes. For any $0<c<1$, at most $J/(-\log c)$ rounds have correct-label probability at most $c$. Fresh posterior-sampled predictions make at most $J$ mistakes in expectation, provided each input is selected before that round's fresh prediction randomness.
\end{corollary}
\begin{proof}
At each finite horizon, the product of successive correct-label probabilities is the prior probability of agreeing with all labels revealed along that sequence. The event $E_y$ is contained in this agreement event, so the product is at least $e^{-J}$. Take logarithms and then increase the horizon. The remaining assertions follow from $1-p\le-\log p$ and the indicated thresholds. The sampling assertion uses conditional expectation under the stated nonanticipation rule.
\end{proof}

The input order can be adaptive and need not be IID. This is a separate exact-sign learner, rather than the noisy Gaussian-likelihood learner of the regression theorems. A continuous function taking both signs cannot have a positive sign margin on the whole connected cube; the sign result instead applies to prescribed sets such as the Boolean cube. The regression neighborhood theorem itself is uniform on the entire cube.

\section{Native compression of the neuron-gated prior}
\label{sec:ard-native-compression}
This section concerns the native hierarchical prior, without conditioning any weights. At each hidden vertex let $S_j^\ell\sim\operatorname{Gamma}(\kappa/N,1)$, with fixed $0<\kappa\le1$. All edge factors and all biases are independent symmetric Student variables of fixed degrees of freedom $\nu>2$. An outgoing weight from vertex $j$ is $S_j^\ell G_{ij}^{\ell+1}$. The first layer has no gate on its input coordinates. The hidden activation is globally $A$-Lipschitz, and $b=|\phi(0)|$. Biases are present at every hidden layer and at the linear readout. Constants below may depend on this fixed activation and scalar-law choice.

\begin{theorem}[A polynomial-depth circuit dictionary]
\label{thm:ard-native-compression}
For $0<\delta,\tau\le1/2$, put
\[
 L=H+1+\log\frac{d+1}{\delta\tau}.
\]
With prior probability at least $1-\delta$, the raw depth-$H$ network is uniformly $\tau$-approximable on $[-1,1]^d$ by a rational-weight network of the same activation and hidden depth, with
\[
 \text{width}\le CL,\qquad
 \text{hidden vertices}\le CHL,\qquad
 \text{bits per coefficient}\le CHL.
\]
The approximating networks belong to a deterministic finite dictionary, independent of learner width $N$, whose cardinality obeys
\[
 \log|\mathcal D_{H,d,\delta,\tau}|
 \le C\big[dH L^2+H^2L^3\big].
 \label{eq:ard-dictionary-cardinality}
\]
These are polynomial bounds jointly in depth, input dimension, and logarithmic failure probability and precision.
\end{theorem}
\begin{proof}
Write $m_1=\E|G|<\infty$. The gamma density implies, uniformly for $a=\kappa/N\in(0,1]$ and $0<\epsilon\le1$,
\begin{align*}
 N\E[S\mathbf1_{S\le\epsilon}]&\le C\epsilon,\\
 N\Pr(S>\epsilon)&\le C\log(C/\epsilon),\\
 N\Pr(S>B)&\le C e^{-B/2}\qquad(B\ge1).
\end{align*}
Indeed $1/\Gamma(a)\le Ca$, integration below $\epsilon$ gives the first bound, integration over $(\epsilon,1)$ and $(1,\infty)$ gives the second, and the exponential gamma tail gives the third.

Delete all outgoing columns whose gates are at most $\epsilon$, keeping the same factors and biases elsewhere. Let $z^{(\epsilon)}$ denote the resulting raw function. The gate $S_j^\ell$ is independent of its own feature $h_j^\ell$, because that feature depends only on earlier gates. Thus the usual path majorants give
\[
 P_1\le b+A m_1(d+1),\qquad
 P_\ell\le b+A m_1+A\kappa m_1P_{\ell-1},
\]
where $P_\ell$ bounds the expected uniform magnitude at any native or truncated neuron. For a fixed $K>1$, these imply $P_\ell\le C(d+1)K^\ell$. Subtract the coupled networks one layer at a time. A retained column propagates a previous feature error with expected total absolute gain at most $\kappa m_1$. The deleted-column contribution has expected absolute gain at most $C\epsilon$. Independence from the preceding features therefore gives
\[
 \E\|z-z^{(\epsilon)}\|_\infty
 \le C(d+1)\epsilon(H+1)K^{H+1}.
\]
Choose
\[
 \epsilon=\frac{\delta\tau}{C(d+1)(H+1)K^{H+1}}
\]
with a sufficiently large constant. The raw truncation error exceeds $\tau/2$ with probability at most $\delta/4$. Also $\log(C/\epsilon)\le CL$.

In one hidden layer, the number of retained gates is binomial with mean at most $C\log(C/\epsilon)\le CL$. Its exponential moment is bounded by that of a Poisson variable with the same mean. Since $\log(H/\delta)\le CL$, a Chernoff bound and union over layers show that every layer retains at most $U=\lceil CL\rceil$ vertices, except on an event of probability at most $\delta/4$. Similarly, the maximum of all $HN$ gates is at most
\[
 B_S=C\log(CH/\delta)
\]
except on an event of probability at most $\delta/8$.

After deletion, only the retained vertices can influence later layers. Their induced subnetwork uses at most
\[
 M=(d+1)U+(H-1)U(U+1)+U+1
       \le C(dL+HL^2)
\]
edge factors and biases. The retained index sets depend only on gates and are independent of every factor and bias. The Student tail bound therefore permits a conditional union bound over these at most $M$ variables. Except on an event of probability at most $\delta/8$, all their magnitudes are at most
\[
 B_G=1+(CM/\delta)^{1/\nu}.
\]
Every actual coefficient of the induced subnetwork is then bounded by $B_*=\max\{1,B_S\}B_G$, and $\log(B_*+1)\le CL$. This includes its readout bias. Empty retained layers are allowed; subsequent retained neurons may then be constants due to their biases.

Relabel retained vertices consecutively in each layer; their original learner indices no longer occur. Round every coefficient to a common dyadic mesh of step $q\le1$. A deterministic layerwise Lipschitz estimate bounds the resulting uniform output error by
\[
 Cq(H+1)(d+U+1)K_q^{H+1},
 \qquad K_q=1+b+A(d+U+1)(B_*+1).
\]
Choosing a dyadic $q$ small enough to make this at most $\tau/2$ requires at most $CHL$ bits per coefficient, since $\log K_q\le CL$. The total failure probability used above is at most $3\delta/4$, and the truncation and rounding errors sum to at most $\tau$.

Finally enumerate all layer widths between zero and $U$, all allowed dyadic coefficients, and all edges in their consecutive dense layer blocks. There are at most $M$ coefficients, each with $CHL$ bits, plus at most $CH\log(U+1)$ bits for the layer widths. This gives \eqref{eq:ard-dictionary-cardinality}. The enumeration contains no original learner-neuron addresses and is independent of $N$.
\end{proof}

\begin{corollary}[An actual circuit-complexity converse]
\label{cor:ard-prior-complexity}
Let $\mathcal X\subseteq[-1,1]^d$ and suppose
\[
 p=\mu_N\{\|z-f\|_{\infty,\mathcal X}\le\tau\}>0,
 \qquad0<\tau\le1/2.
\]
Then $f$ has a uniform $2\tau$ approximation on $\mathcal X$ by a rational-weight network of hidden depth $H$, width at most $CL_p$ and coefficient bit length at most $CHL_p$, where
\[
 L_p=H+1+\log\frac{d+1}{p\tau}.
\]
Its description has at most $C[dH L_p^2+H^2L_p^3]$ bits.
\end{corollary}
\begin{proof}
Apply Theorem~\ref{thm:ard-native-compression} with failure probability $p/2$. Its compression event intersects the specified function neighborhood. The triangle inequality proves the approximation claim, and the remaining conclusions are the dictionary bounds.
\end{proof}

\section{The exact sign prior and online mistakes}
\label{sec:ard-sign-prior}
The circuit-complexity converse also applies to the exact sign-function prior under an all-input margin convention. It does not require an upper bound on raw network outputs.

Fix a nonempty prescribed input set $\mathcal X\subseteq[-1,1]^d$ and a margin $a>0$. Let
\[
 M_a=\{w:|z_w(x)|\ge a\text{ for every }x\in\mathcal X\}.
\]
The unconditioned neuron-gate prior has
\begin{equation}
 \mu_N(M_a)\ge e^{-Q_0},\qquad
 Q_0=C\left[H+\log(d+1)+\log(1+a)+\log(1+a^{-1})\right],
 \label{eq:ard-margin-normalizer}
\end{equation}
uniformly in learner width. To see this, condition the sum of the last-layer gates, which is exactly $\operatorname{Gamma}(\kappa,1)$, to be at most
\[
 \epsilon=\min\{1/4,a/[C(d+1)K^H]\}.
\]
Its negative log probability is at most $C\log(C/\epsilon)$. The last-layer features are independent of these gates, and have expected uniform magnitude at most $C(d+1)K^H$. Thus the readout contribution excluding its bias has conditional expected uniform magnitude at most $a/8$, with a sufficiently large constant. With conditional probability at least $3/4$ its magnitude is at most $a/2$. Independently restrict the readout bias to $[7a/4,9a/4]$. The Student density gives negative log probability at most $C[1+\log(1+a)+\log(1+a^{-1})]$. The resulting function is at least $a$ on the entire input cube, proving \eqref{eq:ard-margin-normalizer}.

\begin{corollary}[The margin-conditioned circuit converse]
\label{cor:ard-margin-converse}
Let $f:\mathcal X\to\{-1,1\}$ have conditional sign prior mass
\[
 p=\mu_N\{f(x)z_w(x)\ge a\text{ for all }x\in\mathcal X\mid M_a\}>0.
\]
Then $f$ is the exact sign of a rational-weight network of hidden depth $H$, width at most $CL_*$ and coefficient bit length at most $CHL_*$, where
\[
 L_*=H+1+\log\frac{d+1}{p\min\{1/2,a/2\}}+Q_0.
\]
For ReLU and Boolean inputs, this gives an ordinary Boolean circuit of size polynomial in $d,H,L_*$, with no exponential dependence on depth.
\end{corollary}
\begin{proof}
The unconditioned joint event has mass at least $p e^{-Q_0}$. Apply native compression at failure probability half this mass and tolerance $\tau=\min\{1/2,a/2\}$. Its good set intersects the joint event, and the compressed function has the same strict sign at every prescribed input. This proves the rational-network statement.

For ReLU on Boolean inputs, evaluate the dyadic network using exact signed integer arithmetic with a common power-of-two denominator. At each layer the required bit length increases by at most the coefficient bit bound plus $O(\log(d+U+1))$, where $U\le CL_*$ is the compressed width. Comparisons with zero implement ReLU exactly. Schoolbook addition and multiplication therefore produce an ordinary Boolean circuit of polynomial size in the number of coefficients, depth and coefficient bit length. All of these resources are polynomial in $d,H,L_*$ by Theorem~\ref{thm:ard-native-compression}.
\end{proof}


\section{The infinite-width limit retains shared neurons}
\label{sec:ard-wide-limit}
The neuron-gated prior also has a precise large-width interpretation. Its contrast with an independent sparse-coordinate prior is that the retained objects are whole neurons, shared by every row that uses them.

\begin{proposition}[A gamma-process circuit limit]
\label{prop:ard-wide-limit}
Fix input dimension, hidden depth, activation and scalar prior parameters. As $N\to\infty$, the raw function laws of the neuron-gated prior converge weakly on $C([-1,1]^d)$. At any fixed gate threshold $\epsilon>0$, the truncated laws converge in total variation to the following finite random network:
\begin{itemize}
\item Each hidden layer has an independent Poisson number of retained neurons, with mean
$\Lambda_\epsilon=\kappa\int_\epsilon^\infty s^{-1}e^{-s}\,ds$.
\item Their gates are independent with density
$\kappa s^{-1}e^{-s}\mathbf1_{s>\epsilon}/\Lambda_\epsilon$.
\item All edge factors and biases between the retained neurons have the original independent Student laws. Each retained gate multiplies the complete outgoing column of its neuron.
\end{itemize}
Thus a retained hidden feature is reused by all of its parents, rather than duplicated into independent branches. The full limiting function law is obtained by letting $\epsilon\downarrow0$. Equivalently each layer has a gamma-process gate cloud of intensity $\kappa s^{-1}e^{-s}\,ds$, with finite total gate mass almost surely.
\end{proposition}
\begin{proof}
For $a=\kappa/N$, the native gate density is
$f_a(s)=s^{a-1}e^{-s}/\Gamma(a)$. On $(\epsilon,\infty)$,
\[
 Nf_a(s)\longrightarrow\kappa s^{-1}e^{-s}
\]
in $L^1$, by dominated convergence. The retained count in a layer is binomial with parameter $q_N=\Pr(S>\epsilon)$, and $Nq_N\to\Lambda_\epsilon$. Binomial-to-Poisson convergence holds in total variation. Conditional on the count, the retained gates are independent with density $f_a\mathbf1_{s>\epsilon}/q_N$, which converges in total variation to the displayed density.

The retained index sets depend only on gates. Therefore all factors and biases of the induced subnetwork remain independent at their native laws, independently of the retained gates. Couple the finitely many layer counts and gates, and then use identical factor arrays whenever the counts agree. This proves total-variation convergence of the truncated network functions. In particular it preserves every use of one retained neuron by several later neurons.

The native truncation estimate in Theorem~\ref{thm:ard-native-compression} is uniform in $N$ and tends to zero with $\epsilon$. Using the bounded metric $\min\{1,\|f-g\|_\infty\}$, the limiting truncated laws are consequently Cauchy in Wasserstein distance, and the same triangle inequality proves convergence of the full finite-width laws.

For the equivalent gate-cloud construction, assign independent Student factors and biases to the countably many atoms in each of the finitely many layers. The gate cloud has total mass $\operatorname{Gamma}(\kappa,1)$ and expected discarded mass $\kappa\int_0^\epsilon e^{-s}ds\le\kappa\epsilon$. The same path majorants give finite expected uniform absolute sums at every row. Thus the recursive infinite sums converge uniformly almost surely, define continuous functions, and have the same truncation limit just obtained.
\end{proof}

\begin{corollary}[The approximation oracle also holds at infinite width]
For any fixed finite data set and bounded Lipschitz final output, the Gaussian-likelihood posterior of the finite-width network converges weakly to the posterior under the limiting function prior. Posterior-mean predictions converge pointwise. Every expected-risk oracle in this paper that is uniform in finite learner width therefore holds for this infinite-width learner as well. The width-independent function-neighborhood and finite-dictionary bounds also pass to the limiting prior.
\end{corollary}
\begin{proof}
For the fixed data set, its likelihood is a bounded, strictly positive continuous functional of the raw function, so normalized likelihood tilting preserves weak convergence. Bounded output evaluations give pointwise convergence of posterior means. The test losses are dominated by $\E(1+|Y|)^2<\infty$, independently of the training sample and width. Dominated convergence therefore passes the expected-risk inequalities to the limit. Closed uniform neighborhoods and the finite union of closed dictionary-approximation neighborhoods pass their probability lower bounds by the closed-set direction of Portmanteau.
\end{proof}

This limiting description is a consequence of the continuous fixed-architecture parameter prior. No discrete width or architecture is sampled by the finite learner. It nevertheless makes the prior's modeling assumption transparent: neuron-level group shrinkage preserves shared computations. An independent sparse-edge prior instead develops a tree-like output ancestry at fixed resolution. The two priors therefore encode different preferences even when their coordinate variances have comparable width scales.


\section{Interpretation and limitations}
The gate is the reason sharing is inexpensive. Each selected neuron incurs one rare order-one gate event, regardless of how many selected neurons use its output. Its possible locations cancel that gate cost. The selected edge factors then pay ordinary finite-dimensional small-ball costs. In contrast, independent rare edges can incur a separate location cost for each shared use.

This construction deliberately imposes continuous group shrinkage. It is not a theorem that every dense Bayesian network has a circuit prior, and it does not establish the same polynomial-depth conclusion for independent Gaussian weights. No empirical claim is made that the prior matches networks trained by gradient descent. The result identifies one explicit, fixed-architecture Bayesian model where a joint width--depth polynomial is possible without a stability restriction on the supported networks.

Student factors are convenient because placing a large coefficient has logarithmic cost. The following alternative uses a standard conditional-Gaussian group-scale model.

\begin{corollary}[Gaussian edge and bias factors]\label{cor:Gaussian-factors}
Replace every Student factor by an independent standard Gaussian, retaining the same gamma gates. The primary learning Theorems~\ref{thm:sharp} and~\ref{thm:main}, and the function-neighborhood upper bounds, remain valid after adding $CPB^2$ to the comparison entropy resource. The native-compression bounds are unchanged. For the margin-conditioned converse at arbitrary margin $a$, replace the normalizer resource $Q_0$ by $Q_0+Ca^2$; for $0<a\le1$ this is absorbed in its constant. For the Boolean restoring implementations, all coefficient caps are numerical constants, so the $\widetilde O((dC+HC^2)/n)$ expected-risk remainder and corresponding online mistake bounds are unchanged.
\end{corollary}
\begin{proof}
A uniform interval of radius $\epsilon$ about any $|a|\le B$ has KL at most $C[\log(C/\epsilon)+B^2]$ relative to a standard Gaussian. All gate and placement calculations are unchanged, as are the moment estimates with $\mu=1$. Summing over $P$ modified factors adds $CPB^2$. The native compression proof needs only a finite first moment and a fixed polynomial tail bound on factors; a Gaussian satisfies both. The Gaussian readout-bias interval used to normalize an arbitrary margin has cost $C[a^2+1+\log(1+a^{-1})]$, which explains the stated adjustment. The Boolean implementations use a fixed coefficient cap and margin, so their additional entropy is absorbed by $CQ$. In Theorem~\ref{thm:scalar-padding}, the Gaussian addition can be bounded by $CP_gB^2+C(H-h+1)$, because the added relay coefficients are numerical constants.
\end{proof}

In this Gaussian variant, outgoing weights share a random scale and are dependent even though their pairwise covariances are zero. The general real-valued comparator bound has polynomial rather than logarithmic coefficient-cap dependence.

The bounded output link is part of the observation model and prediction rule. It allows the statistical argument to compare prediction differences under unbounded sub-Gaussian labels. We do not claim an identical theorem for an unbounded raw-output posterior mean.

\paragraph{Related work and attribution.}
Hierarchical neuron- or group-level shrinkage is an established Bayesian modeling idea. Lee and coauthors~\cite{lee-dependent} study weights of the form $W_{jk}=\sqrt{\lambda_j}V_{jk}$, with Gaussian edge factors and a shared random scale per source neuron, including width-dependent scale arrays, non-Gaussian wide-network limits, sparsity and pruning. Our Gaussian-factor variant is a construction of this shared-node-scale type, with $\lambda_j=(S_j)^2$ and $S_j\sim\operatorname{Gamma}(\kappa/N,1)$. We do not claim to introduce dependent neuron scales, their L\'evy limit interpretation, or pruning as general ideas. Dense independent heavy-tailed network priors also already have posterior-contraction results~\cite{castillo-egels}; neither heavy tails nor deterministic architectures alone are new. The finite-sample statement here is an agnostic comparator oracle, uniform over every prescribed finite learner width and polynomial jointly in comparator width and depth. The sharp expectation proof follows the leave-one-out strategy of H\o gsgaard, Rebeschini and Wegel~\cite{hogsgaard}; the arbitrary-prior and conditional-sub-Gaussian version used here is proved in full.

\paragraph{Verification.}
The comparison proof and its statistical interface have undergone independent mathematical review by AI agents. The accompanying verification record identifies the exact sources. This paper has not been formalized in Lean.
\begin{thebibliography}{9}
\bibitem{castillo-egels} I. Castillo and P. Egels. Posterior and Variational Inference for Deep Neural Networks with Heavy-Tailed Weights. \emph{Journal of Machine Learning Research} 26(122), 1--58, 2025.
\bibitem{lee-dependent} H. Lee, F. Ayed, P. Jung, J. Lee, H. Yang, and F. Caron. Deep Neural Networks with Dependent Weights: Gaussian Process Mixture Limit, Heavy Tails, Sparsity and Compressibility. \emph{Journal of Machine Learning Research} 24(289), 1--78, 2023.
\bibitem{hogsgaard} M. M. H\o gsgaard, P. Rebeschini, and T. Wegel. Aggregation with Exponential Weights is Optimal in Expectation. arXiv:2607.02247, 2026.
\end{thebibliography}
\end{document}
